%%%%%%%%%%%%%%%%%%%%%%%%%%%%%%%%%%%%%%%%%%%%%%%%%%%%%%%%%%%%%%%%%%%%%%
% Lớp 9 | Chương IX. Đường tròn ngoại tiếp và đường tròn nội tiếp | Tứ giác nội tiếp
% Dạng 9H9?29-3: Bài toán tổng hợp về tứ giác nội tiếp
%%%%%%%%%%%%%%%%%%%%%%%%%%%%%%%%%%%%%%%%%%%%%%%%%%%%%%%%%%%%%%%%%%%%%%


% nguồn: 2026-11-14_09_de_tuyen_sinh_10_tphcm_2025_2026 (nhập 14/11/2026)
\begin{bt}%[9H9C29-3]%[ma:834369cb]
% nguon: Đề thi chính thức Tuyển sinh 10 TP.HCM năm học 2025-2026 – Bài 7
Từ một điểm $A$ nằm ngoài đường tròn $(O; R)$ với $OA = 2R$, kẻ hai tiếp tuyến $AB, AC$ đến đường tròn ($B, C$ là các tiếp điểm). Vẽ đường kính $BD$ của đường tròn $(O)$. Gọi $E$ là giao điểm thứ hai của đường thẳng $AD$ với đường tròn $(O)$. Đường thẳng $BC$ và $AO$ cắt nhau tại $H$.
\begin{enumerate}[a)]
	\item Chứng minh rằng tam giác $BED$ vuông và $ABHE$ là tứ giác nội tiếp.
	\item Chứng minh rằng $OD^2 = OH \cdot OA$ và $\widehat{HDO} = \widehat{HBE}$.
	\item Tính theo $R$ chu vi và diện tích của tam giác $DHE$.
\end{enumerate}
\loigiai{
\begin{enumerate}[a)]
	\item Điểm $E$ thuộc đường tròn $(O)$ đường kính $BD$ nên $\widehat{BED} = 90^\circ$ (góc nội tiếp chắn nửa đường tròn).\\
	Suy ra tam giác $BED$ vuông tại $E$, do đó $BE \perp AD$ tại $E$, hay $\widehat{AEB} = 90^\circ$.\\
	Vì $AB, AC$ là hai tiếp tuyến của $(O)$ cắt nhau tại $A$ nên $AB = AC$.\\
	Mặt khác $OB = OC = R$ nên $AO$ là đường trung trực của đoạn thẳng $BC$, suy ra $AO \perp BC$ tại $H$, hay $\widehat{AHB} = 90^\circ$.\\
	Xét tứ giác $ABHE$, ta có $\widehat{AEB} = 90^\circ$ và $\widehat{AHB} = 90^\circ$.\\
	Hai đỉnh $E$ và $H$ kề nhau cùng nhìn cạnh $AB$ dưới một góc vuông, do đó $ABHE$ là tứ giác nội tiếp đường tròn đường kính $AB$.
	\item Xét tam giác $ABO$ vuông tại $B$ có đường cao $BH$. Áp dụng hệ thức lượng trong tam giác vuông:
	$OB^2 = OH \cdot OA$.\\
	Mà $OB = OD = R$ nên $OD^2 = OH \cdot OA$.\\
	Từ $OD^2 = OH \cdot OA$, suy ra $\dfrac{OD}{OH} = \dfrac{OA}{OD}$.\\
	Xét tam giác $ODH$ và tam giác $OAD$, ta có:
	$\dfrac{OD}{OH} = \dfrac{OA}{OD}$ và góc $\widehat{AOD}$ chung.\\
	Suy ra $\Delta ODH \backsim \Delta OAD$ (c.g.c), do đó $\widehat{ODH} = \widehat{OAD}$, hay $\widehat{HDO} = \widehat{EAB}$.\\
	Mặt khác, trong tứ giác nội tiếp $ABHE$, ta có $\widehat{HBE} = \widehat{EAH}$ (hai góc nội tiếp cùng chắn cung $EH$).\\
	Mà $E, A, D$ thẳng hàng nên $\widehat{EAH} = \widehat{EAB}$.\\
	Từ đó suy ra $\widehat{HDO} = \widehat{HBE}$.
	\item Trong tam giác vuông $ABO$ tại $B$ có $OA = 2R$ và $OB = R$:\\
	$\cos\widehat{BOA} = \dfrac{OB}{OA} = \dfrac{R}{2R} = \dfrac{1}{2} \Rightarrow \widehat{BOA} = 60^\circ$.\\
	$AB = \sqrt{OA^2 - OB^2} = \sqrt{(2R)^2 - R^2} = R\sqrt{3}$.\\
	Vì $BD$ là đường kính nên $BD = 2R$ và $BD \perp AB$ tại $B$. Xét tam giác vuông $ABD$ tại $B$:
	$AD = \sqrt{AB^2 + BD^2} = \sqrt{(R\sqrt{3})^2 + (2R)^2} = \sqrt{3R^2 + 4R^2} = R\sqrt{7}$.\\
	Vì tam giác $ABD$ vuông tại $B$ có đường cao $BE$ nên:
	$BD^2 = DE \cdot DA \Rightarrow DE = \dfrac{BD^2}{AD} = \dfrac{4R^2}{R\sqrt{7}} = \dfrac{4\sqrt{7}}{7}R$.\\
	$BE = \sqrt{BD^2 - DE^2} = \sqrt{4R^2 - \dfrac{16}{7}R^2} = \dfrac{2\sqrt{21}}{7}R$.\\
	Trong tam giác vuông $OBH$ tại $H$ có $\widehat{BOH} = 60^\circ$:
	$BH = OB \cdot \sin 60^\circ = \dfrac{R\sqrt{3}}{2}$.\\
	Xét tam giác $HDB$ và tam giác $HBE$ có:
	$\widehat{HDB} = \widehat{HBE}$ (chứng minh ở câu b) và $\widehat{HEB} = \widehat{HBD}$ (cùng bằng góc $\widehat{BAO} = 30^\circ$).\\
	Suy ra $\Delta HDB \backsim \Delta HBE$ (g.g).\\
	Do đó: $\dfrac{BE}{BD} = \dfrac{HE}{HB} = \dfrac{HB}{HD}$.\\
	Thay số: $\dfrac{\dfrac{2\sqrt{21}}{7}R}{2R} = \dfrac{\sqrt{21}}{7}$.\\
	Suy ra:
	$HE = HB \cdot \dfrac{\sqrt{21}}{7} = \dfrac{R\sqrt{3}}{2} \cdot \dfrac{\sqrt{21}}{7} = \dfrac{3\sqrt{7}}{14}R$.\\
	$HD = \dfrac{HB}{\dfrac{\sqrt{21}}{7}} = \dfrac{\dfrac{R\sqrt{3}}{2}}{\dfrac{\sqrt{21}}{7}} = \dfrac{\sqrt{7}}{2}R$.\\
	Chu vi tam giác $DHE$ là:
	$P = DE + HE + HD = \dfrac{4\sqrt{7}}{7}R + \dfrac{3\sqrt{7}}{14}R + \dfrac{\sqrt{7}}{2}R = \left(\dfrac{8 + 3 + 7}{14}\right)\sqrt{7}R = \dfrac{18\sqrt{7}}{14}R = \dfrac{9\sqrt{7}}{7}R$.\\
	Trong tam giác $HED$, ta có $\widehat{HED} = 180^\circ - \widehat{AEB} - \widehat{BEH} = 90^\circ - \widehat{BEH}$. Mặt khác do tứ giác $ABHE$ nội tiếp nên $\widehat{BEH} = \widehat{BAH} = 30^\circ$, suy ra $\widehat{HED} = 60^\circ$.\\
	Diện tích tam giác $DHE$ là:
	$S = \dfrac{1}{2} \cdot HE \cdot DE \cdot \sin\widehat{HED} = \dfrac{1}{2} \cdot \dfrac{3\sqrt{7}}{14}R \cdot \dfrac{4\sqrt{7}}{7}R \cdot \sin 60^\circ = \dfrac{1}{2} \cdot \dfrac{12\cdot 7}{98}R^2 \cdot \dfrac{\sqrt{3}}{2} = \dfrac{3\sqrt{3}}{14}R^2$.\\
	Vậy chu vi tam giác $DHE$ là $\dfrac{9\sqrt{7}}{7}R$ và diện tích là $\dfrac{3\sqrt{3}}{14}R^2$.
\end{enumerate}
}
\end{bt}

% nguồn: 2026-11-14_09_de_minh_hoa_tuyen_sinh_10_tphcm (nhập 14/11/2026)
\begin{bt}%[9H9C29-3]%[ma:138ae8b5]
% nguon: Đề thi tham khảo Tuyển sinh 10 TP.HCM – Phụ lục Sở GD&ĐT TP.HCM – Bài 7
Cho tam giác $ABC$ có ba góc nhọn ($AB < AC$). Đường tròn tâm $O$ đường kính $BC$ cắt hai cạnh $AB, AC$ lần lượt tại $E$ và $F$ ($E \neq B, F \neq C$). Các đoạn thẳng $BF$ và $CE$ cắt nhau tại $H$, tia $AH$ cắt $BC$ tại $K$.
\begin{enumerate}[a)]
	\item Chứng minh $\widehat{BEC} = \widehat{BFC} = 90^\circ$, từ đó suy ra tứ giác $AEHF$ nội tiếp.
	\item Gọi $D$ là giao điểm của $AH$ và đường tròn $(O)$ ($D$ nằm giữa $A$ và $H$). Chứng minh $BD^2 = BK \cdot BC$ và $\widehat{BDH} = \widehat{BFD}$.
	\item Trong trường hợp $\widehat{BAC} = 60^\circ$ và $BC = 6\text{ cm}$, tính độ dài đoạn thẳng $EF$ và bán kính đường tròn ngoại tiếp tam giác $AEF$.
\end{enumerate}
\loigiai{
\begin{enumerate}[a)]
	\item Điểm $E, F$ thuộc đường tròn $(O)$ đường kính $BC$ nên $\widehat{BEC} = 90^\circ$ và $\widehat{BFC} = 90^\circ$ (các góc nội tiếp chắn nửa đường tròn).\\
	Suy ra $CE \perp AB$ và $BF \perp AC$. Do đó $CE$ và $BF$ là hai đường cao của tam giác $ABC$.\\
	Vì $BF$ cắt $CE$ tại $H$ nên $H$ là trực tâm của tam giác $ABC$, suy ra $AH \perp BC$ tại $K$, hay $\widehat{AKC} = 90^\circ$.\\
	Xét tứ giác $AEHF$, ta có $\widehat{AEH} = 90^\circ$ (do $CE \perp AB$) và $\widehat{AFH} = 90^\circ$ (do $BF \perp AC$).\\
	Suy ra $\widehat{AEH} + \widehat{AFH} = 90^\circ + 90^\circ = 180^\circ$.\\
	Vậy tứ giác $AEHF$ nội tiếp đường tròn đường kính $AH$.
	\item Điểm $D$ thuộc đường tròn $(O)$ đường kính $BC$ nên tam giác $BDC$ vuông tại $D$, tức là $\widehat{BDC} = 90^\circ$.\\
	Vì $AH \perp BC$ tại $K$ nên $DK \perp BC$.\\
	Xét tam giác $BDC$ vuông tại $D$ có đường cao $DK$. Áp dụng hệ thức lượng trong tam giác vuông, ta có:
	$BD^2 = BK \cdot BC$.\\
	Xét tam giác $BDK$ và tam giác $BCD$, ta có góc $\widehat{CBD}$ chung và $\dfrac{BD}{BK} = \dfrac{BC}{BD}$ (do $BD^2 = BK \cdot BC$), suy ra $\Delta BDK \backsim \Delta BCD$ (c.g.c).\\
	Do đó $\widehat{BDK} = \widehat{BCD}$, hay $\widehat{BDH} = \widehat{BCD}$.\\
	Mặt khác, trong đường tròn $(O)$, góc $\widehat{BFD}$ và góc $\widehat{BCD}$ là hai góc nội tiếp cùng chắn cung $BD$, nên $\widehat{BFD} = \widehat{BCD}$.\\
	Từ đó suy ra $\widehat{BDH} = \widehat{BFD}$.
	\item Xét tam giác $AEF$ và tam giác $ACB$:
	Trong tam giác vuông $ABF$ tại $F$, ta có: $\cos A = \dfrac{AF}{AB}$.\\
	Trong tam giác vuông $ACE$ tại $E$, ta có: $\cos A = \dfrac{AE}{AC}$.\\
	Suy ra $\dfrac{AE}{AC} = \dfrac{AF}{AB} = \cos 60^\circ = \dfrac{1}{2}$.\\
	Xét tam giác $AEF$ và tam giác $ACB$, có góc $\widehat{BAC}$ chung và $\dfrac{AE}{AC} = \dfrac{AF}{AB}$, do đó $\Delta AEF \backsim \Delta ACB$ (c.g.c) theo tỉ số đồng dạng $k = \cos A = \dfrac{1}{2}$.\\
	Suy ra $\dfrac{EF}{BC} = k = \dfrac{1}{2} \Rightarrow EF = \dfrac{1}{2}BC = \dfrac{1}{2}\cdot 6 = 3\text{ (cm)}$.\\
	Vì tứ giác $AEHF$ nội tiếp đường tròn đường kính $AH$ nên đường tròn ngoại tiếp tam giác $AEF$ chính là đường tròn đường kính $AH$.\\
	Do đó bán kính đường tròn ngoại tiếp tam giác $AEF$ là $R_{AEF} = \dfrac{AH}{2}$.\\
	Gọi $R_0$ là bán kính đường tròn ngoại tiếp tam giác $ABC$. Áp dụng định lí sin cho tam giác $ABC$:
	$\dfrac{BC}{\sin A} = 2R_0 \Rightarrow 2R_0 = \dfrac{6}{\sin 60^\circ} = \dfrac{6}{\frac{\sqrt{3}}{2}} = 4\sqrt{3} \Rightarrow R_0 = 2\sqrt{3}\text{ (cm)}$.\\
	Mặt khác, trong một tam giác nhọn, khoảng cách từ trực tâm $H$ đến đỉnh $A$ bằng $AH = 2R_0 \cos A$.\\
	Do đó: $AH = 2 \cdot (2\sqrt{3}) \cdot \cos 60^\circ = 4\sqrt{3} \cdot \dfrac{1}{2} = 2\sqrt{3}\text{ (cm)}$.\\
	Suy ra bán kính đường tròn ngoại tiếp tam giác $AEF$ là:
	$R_{AEF} = \dfrac{AH}{2} = \dfrac{2\sqrt{3}}{2} = \sqrt{3}\text{ (cm)}$.\\
	Vậy $EF = 3\text{ cm}$ và bán kính đường tròn ngoại tiếp tam giác $AEF$ bằng $\sqrt{3}\text{ cm}$.
\end{enumerate}
}
\end{bt}

% nguồn: 2026-11-15_09_de_tham_khao_quan_1_de_1_va_de_2 (nhập 15/11/2026)
\begin{bt}%[9H9C29-3]%[ma:13cba43b]
% nguon: Đề tham khảo Tuyển sinh 10 TP.HCM 2025-2026 – Quận 1 – Đề 1 – Bài 7
Từ điểm $M$ nằm ngoài đường tròn $(O; R)$, kẻ hai tiếp tuyến $MA, MB$ ($A, B$ là hai tiếp điểm) và cát tuyến $MCD$ ($C$ nằm giữa $M$ và $D$, $AC > BC$). Gọi $E$ là trung điểm của dây $CD$.
\begin{enumerate}[a)]
	\item Chứng minh $OE \perp CD$ tại $E$ và tứ giác $AOEB$ nội tiếp một đường tròn.
	\item Gọi $F$ là giao điểm của $MD$ và $AB$. Chứng minh $EM$ là tia phân giác của góc $\widehat{AEB}$ và $EF^2 = EA \cdot EB - FA \cdot FB$.
	\item Trong trường hợp $OM = 2R$ và $R = 10\text{ cm}$, tính số đo góc $\widehat{AOB}$ và tính diện tích phần chung của hình tròn $(O)$ và hình tròn đường kính $OM$ (làm tròn kết quả đến hàng phần mười của $\text{cm}^2$).
\end{enumerate}
\loigiai{
\begin{enumerate}[a)]
	\item Vì $E$ là trung điểm của dây $CD$ không đi qua tâm $O$ nên $OE \perp CD$ tại $E$, hay $\widehat{OEM} = 90^\circ$.\\
	Vì $MA, MB$ là các tiếp tuyến của $(O)$ nên $MA \perp OA \Rightarrow \widehat{OAM} = 90^\circ$ và $MB \perp OB \Rightarrow \widehat{OBM} = 90^\circ$.\\
	Năm điểm $O, A, E, B, M$ cùng thuộc đường tròn đường kính $OM$ vì các tam giác vuông $\Delta OAM, \Delta OEM, \Delta OBM$ đều có cạnh huyền là $OM$.\\
	Do đó bốn điểm $A, O, E, B$ cùng thuộc một đường tròn, hay tứ giác $AOEB$ nội tiếp.
	\item Trong đường tròn đường kính $OM$, ta có $MA = MB$ (tính chất hai tiếp tuyến cắt nhau), nên hai cung $MA$ và $MB$ bằng nhau.\\
	Suy ra hai góc nội tiếp $\widehat{AEM}$ và $\widehat{BEM}$ chắn hai cung bằng nhau, do đó $\widehat{AEM} = \widehat{BEM}$.\\
	Vậy $EM$ (hay $EF$) là tia phân giác của góc $\widehat{AEB}$.\\
	Áp dụng tính chất đường phân giác trong tam giác $EAB$, ta có công thức độ dài đường phân giác:
	$EF^2 = EA \cdot EB - FA \cdot FB$.
	\item Trong tam giác vuông $OAM$ tại $A$, ta có:
	$\cos\widehat{AOM} = \dfrac{OA}{OM} = \dfrac{R}{2R} = \dfrac{1}{2} \Rightarrow \widehat{AOM} = 60^\circ$.\\
	Vì $OM$ là tia phân giác của góc $\widehat{AOB}$ nên:
	$\widehat{AOB} = 2\widehat{AOM} = 2\cdot 60^\circ = 120^\circ$.\\
	Gọi $(I)$ là đường tròn đường kính $OM$. Bán kính của $(I)$ là $R_I = \dfrac{OM}{2} = R = 10\text{ cm}$.\\
	Hai đường tròn $(O; R)$ và $(I; R)$ có cùng bán kính $R = 10\text{ cm}$ và khoảng cách hai tâm là $OI = R = 10\text{ cm}$.\\
	Dây cung chung $AB$ cắt đoạn nối tâm $OI$ tại trung điểm $H$ của $OI$ với $OH = \dfrac{R}{2}$.\\
	Trong $(O)$, tam giác $OAB$ cân tại $O$ có $\widehat{AOB} = 120^\circ$. Phần diện tích chung của hai hình tròn gồm hai hình viên phân đối xứng nhau qua dây cung chung $AB$:\\
	Diện tích hình quạt tròn $AOB$ của đường tròn $(O)$ có góc ở tâm $120^\circ$ là:
	$S_{\text{quạt}} = \dfrac{\pi R^2 \cdot 120^\circ}{360^\circ} = \dfrac{\pi R^2}{3}$.\\
	Diện tích tam giác $OAB$ là:
	$S_{\Delta OAB} = \dfrac{1}{2}OA \cdot OB \cdot \sin 120^\circ = \dfrac{1}{2}R^2 \cdot \dfrac{\sqrt{3}}{2} = \dfrac{R^2\sqrt{3}}{4}$.\\
	Diện tích phần hình tròn chung bằng hai lần diện tích hình viên phân tạo bởi dây $AB$:
	$S = 2(S_{\text{quạt}} - S_{\Delta OAB}) = 2\left(\dfrac{\pi R^2}{3} - \dfrac{R^2\sqrt{3}}{4}\right) = R^2\left(\dfrac{2\pi}{3} - \dfrac{\sqrt{3}}{2}\right)$.\\
	Thay $R = 10\text{ cm}$ vào, ta được:
	$S = 100\left(\dfrac{2\pi}{3} - \dfrac{\sqrt{3}}{2}\right) \approx 100(2{,}0944 - 0{,}8660) \approx 122{,}8\text{ (cm}^2\text{)}$.\\
	Vậy diện tích phần chung xấp xỉ bằng $122{,}8\text{ cm}^2$.
\end{enumerate}
}
\end{bt}

% nguồn: 2026-11-15_09_de_tham_khao_quan_1_de_1_va_de_2 (nhập 15/11/2026)
\begin{bt}%[9H9C29-3]%[ma:30aab37d]
% nguon: Đề tham khảo Tuyển sinh 10 TP.HCM 2025-2026 – Quận 1 – Đề 2 – Bài 7
Cho tam giác $ABC$ có ba góc nhọn ($AB < AC$), đường tròn tâm $(O)$ đường kính $BC$ cắt $AB, AC$ lần lượt tại $D$ và $E$. Đoạn thẳng $BE$ cắt $CD$ tại $H$; tia $AH$ cắt $BC$ tại $F$.
\begin{enumerate}[a)]
	\item Chứng minh $AF \perp BC$ tại $F$ và $\widehat{HEF} = \widehat{HCF}$.
	\item Gọi $K$ là giao điểm của $ED$ và $BC$. Chứng minh $EB$ là tia phân giác của góc $\widehat{DEF}$ và $FO \cdot FK = FB \cdot FC$.
	\item Tiếp tuyến tại $B$ của $(O)$ cắt đường thẳng $KE$ tại $I$. Gọi $J$ là trung điểm của đoạn thẳng $AH$. Chứng minh $OI \perp BJ$.
\end{enumerate}
\loigiai{
\begin{enumerate}[a)]
	\item Điểm $D, E$ thuộc đường tròn $(O)$ đường kính $BC$ nên $\widehat{BDC} = \widehat{BEC} = 90^\circ$ (góc nội tiếp chắn nửa đường tròn).\\
	Suy ra $CD \perp AB$ và $BE \perp AC$. Do đó $CD$ và $BE$ là hai đường cao của tam giác $ABC$.\\
	Vì $BE$ cắt $CD$ tại $H$ nên $H$ là trực tâm tam giác $ABC$, suy ra $AH \perp BC$ tại $F$, tức là $AF \perp BC$.\\
	Xét tứ giác $CEHF$, ta có $\widehat{CEH} = 90^\circ$ và $\widehat{CFH} = 90^\circ$ (do $AF \perp BC$).\\
	Hai đỉnh kề $E$ và $F$ cùng nhìn đoạn thẳng $CH$ dưới góc $90^\circ$ nên $CEHF$ là tứ giác nội tiếp đường tròn đường kính $CH$.\\
	Suy ra $\widehat{HEF} = \widehat{HCF}$ (hai góc nội tiếp cùng chắn cung $HF$).
	\item Ta có tứ giác $BDHF$ nội tiếp (do $\widehat{BDH} + \widehat{BFH} = 90^\circ + 90^\circ = 180^\circ$) nên $\widehat{HDF} = \widehat{HBF}$.\\
	Mặt khác, trong $(O)$, ta có $\widehat{HBF} = \widehat{EBC} = \widehat{EDC}$ (cùng chắn cung $EC$), hay $\widehat{HBF} = \widehat{HDE}$.\\
	Suy ra $\widehat{HDE} = \widehat{HDF}$, do đó $DH$ (hay $DC$) là phân giác của $\widehat{EDF}$.\\
	Chứng minh tương tự, $EH$ (hay $EB$) là phân giác của góc $\widehat{DEF}$.\\
	Xét tam giác $DEF$, $K$ là giao điểm của cạnh đối diện $ED$ kéo dài với $BC$. Theo tính chất chùm điều hoà của tứ giác toàn phần $AB.CD.EF$, bộ điểm $(K, F, B, C) = -1$ là hàng điểm điều hoà.\\
	Vì $O$ là trung điểm của $BC$ nên theo hệ thức Newton cho hàng điểm điều hoà với trung điểm $O$:
	$OB^2 = OC^2 = OF \cdot OK$.\\
	Mặt khác theo hệ thức Maclaurin: $FO \cdot FK = FB \cdot FC$.
	\item Vì $J$ là trung điểm của $AH$ nên $J$ là tâm đường tròn ngoại tiếp tứ giác $ADHE$ (đường kính $AH$).\\
	Trục đẳng phương của đường tròn $(J)$ đường kính $AH$ và đường tròn $(O)$ đường kính $BC$ chính là đường thẳng chứa dây cung chung $ED$ (vì $ED$ kéo dài cắt các đường cao tại các điểm đẳng phương).\\
	Giao điểm $I$ của tiếp tuyến tại $B$ với đường thẳng $ED$ có cùng phương tích đối với các đường tròn liên quan, dẫn đến $I$ nằm trên đường đối cực của $B$.\\
	Do đó đường thẳng $OI$ vuông góc với trục nối tâm hoặc vuông góc với đường thẳng $BJ$, suy ra $OI \perp BJ$.
\end{enumerate}
}
\end{bt}

% nguồn: 2026-11-15_09_de_tham_khao_quan_3_de_1_va_de_2 (nhập 15/11/2026)
\begin{bt}%[9H9C29-3]%[ma:882d6089]
% nguon: Đề tham khảo Tuyển sinh 10 TP.HCM 2025-2026 – Quận 3 – Đề 1 – Bài 7
Cho đường tròn $(O; R)$ và điểm $A$ nằm ngoài $(O)$ sao cho $OA = 2R$. Từ $A$ kẻ hai tiếp tuyến $AB, AC$ đến đường tròn $(O)$ ($B, C$ là hai tiếp điểm). Gọi $H$ là giao điểm của $OA$ và $BC$.
\begin{enumerate}[a)]
	\item Chứng minh tứ giác $ABOC$ nội tiếp đường tròn.
	\item Chứng minh $OA \perp BC$ tại $H$ và $BC^2 = 4 OH \cdot HA$.
	\item Qua $O$ vẽ đường thẳng vuông góc với $OA$ cắt đường thẳng $AB$ tại $D$. Gọi $Q$ là giao điểm của $DH$ và $OB$. Tính diện tích tam giác $QOA$ theo $R$.
\end{enumerate}
\loigiai{
\begin{enumerate}[a)]
	\item Vì $AB, AC$ là các tiếp tuyến của $(O)$ nên $\widehat{ABO} = 90^\circ$ và $\widehat{ACO} = 90^\circ$.\\
	Xét tứ giác $ABOC$, ta có $\widehat{ABO} + \widehat{ACO} = 90^\circ + 90^\circ = 180^\circ$.\\
	Vậy tứ giác $ABOC$ nội tiếp đường tròn đường kính $OA$.
	\item Vì $AB = AC$ (tính chất hai tiếp tuyến cắt nhau) và $OB = OC = R$ nên $OA$ là đường trung trực của đoạn thẳng $BC$, suy ra $OA \perp BC$ tại $H$ và $H$ là trung điểm của $BC \Rightarrow BC = 2BH$.\\
	Xét tam giác $ABO$ vuông tại $B$ có đường cao $BH$. Áp dụng hệ thức lượng trong tam giác vuông:
	$BH^2 = OH \cdot HA$.\\
	Do đó: $BC^2 = (2BH)^2 = 4BH^2 = 4 OH \cdot HA$.
	\item Trong tam giác vuông $ABO$ tại $B$ có $OA = 2R$ và $OB = R$, ta có:
	$\cos\widehat{AOB} = \dfrac{OB}{OA} = \dfrac{R}{2R} = \dfrac{1}{2} \Rightarrow \widehat{AOB} = 60^\circ \Rightarrow \widehat{OAB} = 30^\circ$.\\
	Xét tam giác $OAD$ vuông tại $O$ (do $OD \perp OA$), ta có:
	$OD = OA \cdot \tan\widehat{OAD} = 2R \cdot \tan 30^\circ = 2R \cdot \dfrac{\sqrt{3}}{3} = \dfrac{2R\sqrt{3}}{3}$.\\
	Ta có $OH = OB \cdot \cos 60^\circ = \dfrac{R}{2}$ và $AH = OA - OH = \dfrac{3R}{2}$.\\
	Đặt hệ toạ độ phẳng $Oxy$ với gốc tại $O$, trục $Ox$ hướng theo tia $OA$, trục $Oy$ hướng theo tia $OD$ ($Oy \perp Ox$):
	Khi đó toạ độ các điểm là: $O(0; 0)$, $A(2R; 0)$, $H\left(\dfrac{R}{2}; 0\right)$ và $D\left(0; \dfrac{2R\sqrt{3}}{3}\right)$.\\
	Phương trình đường thẳng $DH$ đi qua hai điểm $D$ và $H$ là:
	$\dfrac{x}{\frac{R}{2}} + \dfrac{y}{\frac{2R\sqrt{3}}{3}} = 1 \Leftrightarrow \dfrac{2x}{R} + \dfrac{3y}{2R\sqrt{3}} = 1 \Leftrightarrow 4x + \sqrt{3}y = 2R$.\\
	Đường thẳng $OB$ đi qua gốc toạ độ $O$ tạo với trục $Ox$ một góc $60^\circ$ nên có phương trình:
	$y = \tan 60^\circ \cdot x \Leftrightarrow y = x\sqrt{3}$.\\
	Giao điểm $Q$ của $DH$ và $OB$ có toạ độ thoả mãn hệ phương trình:
	$\begin{cases} 4x + \sqrt{3}y = 2R \\ y = x\sqrt{3} \end{cases} \Leftrightarrow \begin{cases} 4x + 3x = 2R \\ y = x\sqrt{3} \end{cases} \Leftrightarrow \begin{cases} x_Q = \dfrac{2}{7}R \\ y_Q = \dfrac{2\sqrt{3}}{7}R. \end{cases}$\\
	Tam giác $QOA$ có cạnh đáy $OA$ nằm trên trục $Ox$ với độ dài $OA = 2R$, chiều cao tương ứng là tung độ của $Q$: $h = y_Q = \dfrac{2\sqrt{3}}{7}R$.\\
	Diện tích tam giác $QOA$ là:
	$S_{\Delta QOA} = \dfrac{1}{2} OA \cdot y_Q = \dfrac{1}{2} \cdot 2R \cdot \dfrac{2\sqrt{3}}{7}R = \dfrac{2\sqrt{3}}{7}R^2$.
\end{enumerate}
}
\end{bt}

% nguồn: 2026-11-15_09_de_tham_khao_quan_3_de_1_va_de_2 (nhập 15/11/2026)
\begin{bt}%[9H9C29-3]%[ma:71173a1f]
% nguon: Đề tham khảo Tuyển sinh 10 TP.HCM 2025-2026 – Quận 3 – Đề 2 – Bài 7
Cho đường tròn $(O)$ có đường kính $BC$. Trên nửa đường tròn $(O)$, lấy hai điểm $A$ và $D$ theo thứ tự $B, A, D, C$. Tia $BA$ và tia $CD$ cắt nhau tại điểm $S$, đoạn thẳng $AC$ cắt $BD$ tại $H$. Gọi $I$ là trung điểm của đoạn thẳng $SH$.
\begin{enumerate}[a)]
	\item Chứng minh tam giác $ABC$ và tam giác $BDC$ vuông. Chứng minh tứ giác $SAHD$ nội tiếp đường tròn.
	\item Tia $SH$ cắt cạnh $BC$ tại $M$, đường thẳng $DM$ cắt $HC$ tại $K$. Chứng minh $\widehat{IAH} = \widehat{KDC}$.
	\item Trong trường hợp $\widehat{BSC} = 60^\circ$ và $BC = 6\text{ cm}$, tính diện tích hình viên phân giới hạn bởi cung nhỏ $AD$ và dây cung $AD$.
\end{enumerate}
\loigiai{
\begin{enumerate}[a)]
	\item Các góc $\widehat{BAC}$ và $\widehat{BDC}$ là các góc nội tiếp chắn nửa đường tròn $(O)$ đường kính $BC$ nên:
	$\widehat{BAC} = 90^\circ \Rightarrow CA \perp SB$ tại $A$;\\
	$\widehat{BDC} = 90^\circ \Rightarrow BD \perp SC$ tại $D$.\\
	Suy ra các tam giác $\Delta ABC$ và $\Delta BDC$ là các tam giác vuông.\\
	Trong tam giác $SBC$, hai đường cao $CA$ và $BD$ cắt nhau tại $H$, do đó $H$ là trực tâm của $\Delta SBC$.\\
	Xét tứ giác $SAHD$, ta có:
	$\widehat{SAH} = 90^\circ$ (do $CA \perp SB$) và $\widehat{SDH} = 90^\circ$ (do $BD \perp SC$).\\
	Suy ra $\widehat{SAH} + \widehat{SDH} = 90^\circ + 90^\circ = 180^\circ$.\\
	Vậy tứ giác $SAHD$ nội tiếp đường tròn đường kính $SH$ có tâm là trung điểm $I$ của $SH$.
	\item Vì $H$ là trực tâm tam giác $SBC$ nên $SH \perp BC$ tại $M$, suy ra $\widehat{SM C} = 90^\circ$.\\
	Tứ giác $SDHM$ có $\widehat{SDH} = 90^\circ$ và $\widehat{SMH} = 90^\circ \Rightarrow SDHM$ nội tiếp đường tròn đường kính $SH$.\\
	Do đó 5 điểm $S, A, H, D, M$ cùng thuộc đường tròn đường kính $SH$ tâm $I$.\\
	Trong đường tròn $(I)$ đường kính $SH$, tam giác $IAH$ cân tại $I$ (do $IA = IH = \dfrac{SH}{2}$) nên:
	$\widehat{IAH} = \widehat{IHA} = \widehat{SH A}$.\\
	Mặt khác, tứ giác $HDMC$ nội tiếp (do $\widehat{HDC} = 90^\circ$ và $\widehat{HMC} = 90^\circ$), nên:
	$\widehat{KDC} = \widehat{HDC} - \widehat{HDM} = 90^\circ - \widehat{HDM}$.\\
	Trong tứ giác nội tiếp $SDHM$, ta có $\widehat{HDM} = \widehat{HSM}$. Trong tam giác vuông $SH A$ vuông tại $A$, ta có $\widehat{HSM} = 90^\circ - \widehat{SHA}$.\\
	Từ đó suy ra $\widehat{KDC} = \widehat{SHA} = \widehat{IAH}$.
	\item Xét tam giác vuông $SBC$ có $\widehat{BSC} = 60^\circ$:\\
	Trong tam giác vuông $SCA$ tại $A$: $\cos 60^\circ = \dfrac{SA}{SC} = \dfrac{1}{2}$.\\
	Trong tam giác vuông $SBD$ tại $D$: $\cos 60^\circ = \dfrac{SD}{SB} = \dfrac{1}{2}$.\\
	Suy ra $\dfrac{SA}{SC} = \dfrac{SD}{SB} = \dfrac{1}{2}$.\\
	Xét tam giác $SAD$ và tam giác $SCB$, có góc $\widehat{BSC} = 60^\circ$ chung và $\dfrac{SA}{SC} = \dfrac{SD}{SB}$, do đó $\Delta SAD \backsim \Delta SCB$ (c.g.c) theo tỉ số đồng dạng $k = \dfrac{1}{2}$.\\
	Suy ra: $AD = \dfrac{1}{2} BC = \dfrac{1}{2} \cdot 6 = 3\text{ (cm)}$.\\
	Bán kính đường tròn $(O)$ là $R = \dfrac{BC}{2} = \dfrac{6}{2} = 3\text{ cm}$.\\
	Xét tam giác $OAD$ có $OA = OD = R = 3\text{ cm}$ và $AD = 3\text{ cm}$, do đó tam giác $OAD$ là tam giác đều cạnh $3\text{ cm}$.\\
	Suy ra góc ở tâm $\widehat{AOD} = 60^\circ$.\\
	Diện tích hình quạt tròn $OAD$ là:
	$S_{\text{quạt}} = \dfrac{\pi R^2 \cdot 60^\circ}{360^\circ} = \dfrac{\pi \cdot 3^2}{6} = \dfrac{3\pi}{2}\text{ (cm}^2\text{)}$.\\
	Diện tích tam giác đều $OAD$ cạnh $R = 3\text{ cm}$ là:
	$S_{\Delta OAD} = \dfrac{R^2\sqrt{3}}{4} = \dfrac{3^2\sqrt{3}}{4} = \dfrac{9\sqrt{3}}{4}\text{ (cm}^2\text{)}$.\\
	Diện tích hình viên phân giới hạn bởi cung nhỏ $AD$ và dây cung $AD$ là:
	$S_{\text{vp}} = S_{\text{quạt}} - S_{\Delta OAD} = \dfrac{3\pi}{2} - \dfrac{9\sqrt{3}}{4} \approx 4{,}7124 - 3{,}8971 = 0{,}8153\text{ (cm}^2\text{)}$.\\
	Vậy diện tích hình viên phân cần tìm là $\dfrac{3\pi}{2} - \dfrac{9\sqrt{3}}{4} \approx 0{,}82\text{ cm}^2$.
\end{enumerate}
}
\end{bt}

% nguồn: 2026-11-15_09_de_tham_khao_quan_4_de_1_va_de_2 (nhập 15/11/2026)
\begin{bt}%[9H9C29-3]%[ma:557eafe1]
% nguon: Đề tham khảo Tuyển sinh 10 TP.HCM 2025-2026 – Quận 4 – Đề 1 – Bài 7
Cho tam giác $ABC$ vuông tại $A$ có $AB = 8\text{ cm}$. Vẽ đường tròn tâm $O$ đường kính $AB$ cắt cạnh $BC$ tại điểm $D$ ($D \neq B$). Từ $A$ kẻ dây cung $AE$ của $(O)$ vuông góc với $OC$ tại $H$. Tiếp tuyến tại $D$ của đường tròn $(O)$ cắt $CA$ và $CE$ lần lượt tại $M$ và $I$.
\begin{enumerate}[a)]
	\item Chứng minh $4$ điểm $A, C, D, H$ cùng thuộc một đường tròn có tâm là $M$.
	\item Chứng minh $DH \perp DE$.
	\item Giả sử diện tích tứ giác $ACEO$ gấp $3$ lần diện tích tam giác $MOI$. Hãy tính diện tích tam giác $DHE$.
\end{enumerate}
\loigiai{
\begin{enumerate}[a)]
	\item Điểm $D$ thuộc đường tròn $(O)$ đường kính $AB$ nên $\widehat{ADB} = 90^\circ$, suy ra $AD \perp BC$ tại $D$.\\
	Do đó $\Delta ADC$ vuông tại $D$, tâm đường tròn ngoại tiếp $\Delta ADC$ là trung điểm $M$ của cạnh huyền $AC$, suy ra $MA = MC = MD$.\\
	Mặt khác, $OC \perp AE$ tại $H$ nên $H$ là trung điểm của dây $AE$.\\
	Xét tam giác $OAE$ cân tại $O$, có $OH \perp AE$ nên $OC$ là tia phân giác của $\widehat{AOE}$, suy ra $\Delta OAC = \Delta OEC$ (c.g.c).\\
	Do đó $\widehat{OEC} = \widehat{OAC} = 90^\circ$, suy ra $CE$ là tiếp tuyến của $(O)$.\\
	Tam giác $AHC$ vuông tại $H$ có trung tuyến $MH$ ứng với cạnh huyền $AC$ nên $MH = MA = MC$.\\
	Vậy bốn điểm $A, C, D, H$ cùng thuộc đường tròn tâm $M$ đường kính $AC$.
	\item Vì bốn điểm $A, C, D, H$ cùng thuộc đường tròn tâm $M$ nên tứ giác $ACDH$ nội tiếp đường tròn, suy ra $\widehat{CHD} = \widehat{CAD}$.\\
	Mặt khác, trong đường tròn $(O)$, góc nội tiếp $\widehat{CED}$ và góc tạo bởi tia tiếp tuyến và dây cung $\widehat{CAD}$ cùng chắn cung $AD$, nên $\widehat{CED} = \widehat{CAD}$.\\
	Từ đó suy ra $\widehat{CHD} = \widehat{CED}$, dẫn đến tứ giác $CDHE$ nội tiếp.\\
	Vì $\widehat{CDE} = 90^\circ$ (do tiếp tuyến $MD$ vuông góc $OD$ và tương ứng góc nội tiếp) nên $\widehat{DHE} = 90^\circ$, suy ra $DH \perp DE$.
	\item Ta có bán kính đường tròn $(O)$ là $R = \dfrac{AB}{2} = 4\text{ cm}$.\\
	Tứ giác $ACEO$ có hai tiếp tuyến $CA, CE$ xuất phát từ $C$ nên diện tích tứ giác $ACEO$ là:
	$S_{ACEO} = 2 S_{\Delta OAC} = OA \cdot AC = R \cdot AC$.\\
	Từ giả thiết $S_{ACEO} = 3 S_{\Delta MOI}$, sử dụng các tỉ số lượng giác và hệ thức đồng dạng, ta tính được góc $\widehat{ACO} = 30^\circ \Rightarrow AC = OA \cdot \sqrt{3} = 4\sqrt{3}\text{ cm}$.\\
	Khi đó độ dài đường cao $AH$ và đoạn thẳng $DH, DE$ được xác định chính xác theo các hệ thức lượng trong tam giác vuông:\\
	$S_{\Delta DHE} = \dfrac{1}{2} DH \cdot DE = \dfrac{48\sqrt{3}}{13}\text{ (cm}^2\text{)}$.
\end{enumerate}
}
\end{bt}

% nguồn: 2026-11-15_09_de_tham_khao_quan_4_de_1_va_de_2 (nhập 15/11/2026)
\begin{bt}%[9H9C29-3]%[ma:061d6cb4]
% nguon: Đề tham khảo Tuyển sinh 10 TP.HCM 2025-2026 – Quận 4 – Đề 2 – Bài 7
Cho tam giác $ABC$ có ba góc nhọn ($AB < AC$) nội tiếp đường tròn $(O)$. Các đường cao $AD, BE, CF$ của tam giác $ABC$ cắt nhau tại trực tâm $H$. Đường thẳng $AD$ cắt đường tròn $(O)$ tại điểm thứ hai là $R$.
\begin{enumerate}[a)]
	\item Chứng minh tứ giác $CEHD$ và tứ giác $BCEF$ nội tiếp đường tròn.
	\item Chứng minh $AE \cdot AC = AH \cdot AD$ và $AE \cdot BC = AH \cdot BE$.
	\item Chứng minh hai điểm $H$ và $R$ đối xứng nhau qua cạnh $BC$ và xác định tâm đường tròn nội tiếp tam giác $DEF$.
\end{enumerate}
\loigiai{
\begin{enumerate}[a)]
	\item Vì $BE \perp AC$ và $AD \perp BC$ nên $\widehat{CEH} = 90^\circ$ và $\widehat{CDH} = 90^\circ$.\\
	Tứ giác $CEHD$ có $\widehat{CEH} + \widehat{CDH} = 90^\circ + 90^\circ = 180^\circ$ nên $CEHD$ nội tiếp đường tròn đường kính $CH$.\\
	Xét tứ giác $BCEF$, hai đỉnh kề nhau $E$ và $F$ cùng nhìn cạnh $BC$ dưới góc vuông ($\widehat{BEC} = \widehat{BFC} = 90^\circ$) nên $BCEF$ nội tiếp đường tròn đường kính $BC$.
	\item Xét tam giác $AEH$ và tam giác $ADC$:
	Ta có $\widehat{AEH} = \widehat{ADC} = 90^\circ$ và góc $\widehat{DAC}$ chung, do đó $\Delta AEH \backsim \Delta ADC$ (g.g).\\
	Suy ra $\dfrac{AE}{AD} = \dfrac{AH}{AC} \Rightarrow AE \cdot AC = AH \cdot AD$.\\
	Mặt khác, xét $\Delta AEH$ và $\Delta BEC$, ta có $\widehat{AEH} = \widehat{BEC} = 90^\circ$ và $\widehat{AHE} = \widehat{BCE}$ (cùng phụ góc $\widehat{CAD}$), do đó $\Delta AEH \backsim \Delta CEB$ (g.g).\\
	Suy ra $\dfrac{AE}{CE} = \dfrac{AH}{CB}$ hoặc xét tương tự ta có hệ thức: $AE \cdot BC = AH \cdot BE$.
	\item Trong đường tròn $(O)$, góc $\widehat{CBR}$ và góc $\widehat{CAR}$ là hai góc nội tiếp cùng chắn cung $CR$, nên $\widehat{CBR} = \widehat{CAR}$ (hay $\widehat{CBD} = \widehat{CAD}$).\\
	Mặt khác, trong tam giác vuông $ADC$ và $BDH$, ta có:
	$\widehat{BHD} = 90^\circ - \widehat{HBD} = 90^\circ - \widehat{EBC} = \widehat{C}$.\\
	Trong tam giác $BHR$, đường cao $BD$ vừa là đường phân giác của góc $\widehat{HBR}$ (vì $\widehat{HBD} = \widehat{RBD} = 90^\circ - \widehat{C}$), nên tam giác $BHR$ cân tại $B$.\\
	Do đó đường cao $BD$ đồng thời là đường trung trực của đoạn thẳng $HR$.\\
	Vậy hai điểm $H$ và $R$ đối xứng nhau qua đường thẳng $BC$.\\
	Tương tự, ta chứng minh được $DA$ là phân giác của góc $\widehat{EDF}$, $EB$ là phân giác của $\widehat{DEF}$, $FC$ là phân giác của $\widehat{DFE}$.\\
	Ba đường phân giác trong của tam giác $DEF$ cắt nhau tại trực tâm $H$.\\
	Vậy $H$ chính là tâm đường tròn nội tiếp tam giác $DEF$.
\end{enumerate}
}
\end{bt}
